\begin{frame}{자주 묻는 질문 (14.1)}
  \small
  \begin{itemize}
    \item \textbf{Q. \texttt{qpos}\,·\,\texttt{qvel}\,·\,\texttt{ctrl}은?}
      $q_{pos}$ = 위치 자유도(관절 각도, 물체 위치),
      $q_{vel}$ = 그 시간 미분(각속도, 속도) --- MDP 상태
      $\mathcal{S}$가 바로 이 $(q_{pos}, q_{vel})$의 조합.
      \texttt{ctrl} = 액추에이터에 가하는 제어 입력(토크) ---
      MDP의 행동 $\mathcal{A}$에 해당.
    \item \textbf{Q. 왜 굳이 MuJoCo API를 직접 쓰나?}
      기존 환경(Reacher, Ant)은 Gymnasium 래퍼가 편리하지만,
      ``세상에 없는 새로운 로봇''을 시뮬레이션하려면 XML로 직접
      정의해야 한다 --- Gymnasium의 MuJoCo 환경도 내부적으로
      정확히 이 API를 감싸고 있을 뿐이다.
  \end{itemize}
\end{frame}
